\chapter{Introduction}

The project, titled \textit{Biometric-Based Encrypted Vault Implementation}, presents a small secure storage system that combines biometric authentication with encrypted storage management. The main goal is to create a vault that behaves like a normal mounted storage area after it is unlocked. Access to the vault is controlled by a fingerprint sensor and an embedded firmware component. The project uses standard Linux storage encryption mechanisms and focuses on the surrounding system, including enrollment, unlock, lock, formatting, device management, and communication between the host computer and the biometric hardware.

The hardware side of the project is based on an AS608 fingerprint sensor, an STM32 Blue Pill microcontroller, and a PL2303 USB-UART bridge. Fingerprint templates are kept inside the fingerprint sensor, so the host computer does not store raw biometric data. The cryptographic key is managed by the firmware side and becomes available only after a successful fingerprint verification. The key is not permanently stored on the host computer. It is used for the vault operation and then cleared from memory. The transfer between the firmware and the host uses a protected channel, so the key is not exposed as plaintext during communication.

On the storage side, the project uses a Linux Unified Key Setup (LUKS2)/dm-crypt based design with AES-256-XTS encryption. A vault can be created on a physical block device or inside a virtual container backed by a sparse image and loop device. During setup, the system formats the encrypted vault partition, prepares the filesystem, and registers the vault information used by the graphical interface. During normal use, a successful biometric verification unlocks the encrypted storage and mounts it as an accessible filesystem. The user can then work with files through the operating system as usual. When the vault is locked, the mount is removed and the encrypted mapper is closed.

The graphical interface provides the main workflows of the system: creating a vault, formatting storage, enrolling authorized fingerprints during vault creation, unlocking and locking the vault, viewing vault details, and monitoring security-related events. The project also includes safety-oriented behavior such as automatic locking after inactivity, emergency locking when the sensor connection is lost or when the vault's block device is removed, and a force-lock option for cases where files inside the vault are still being used by another process.

The rest of this report explains the project in more detail. Chapter 2 introduces the background technologies used in the system. Chapter 3 describes the system architecture and design decisions. Chapter 4 presents the implementation of the hardware interface, vault manager, and graphical application. Chapter 5 discusses testing and evaluation. Chapter 6 presents the conclusion, limitations, and possible future improvements.
